\begin{papercard}{mcdonald2019}
\cf{Stages} \stagetag{C} \stagetag{I} (also \stagetag{R}). A two-part study of how HCI/CSCW researchers handle reliability in qualitative work. Part 1 is a content analysis of 261 qualitative CHI/CSCW papers (2016--2018). Part 2 gives guidelines on \emph{when} to use agreement or inter-rater reliability (IRR) and how to report it. It is the justification source for \emph{whether} the thesis needs IRR at all.

\cf{Definitions} IRR is a statistical measure of agreement between two or more coders. It merges a positivist quantitative method with interpretivist qualitative methods, and guidelines exist for \emph{how} but rarely for \emph{when} \pc{mcdonald2019}{2}{IRR is a statistical measure of agreement between two or more coders of data}. Reliability concerns reproducibility; validity concerns correctness. Krippendorff's three types: replicability across coders, stability of one coder over time, and accuracy against a standard \pc{mcdonald2019}{4}{Krippendorff proposed three types of reliability: replicability of results across coders}. Calculating IRR yields a claim about the consensus process and not a result \pc{mcdonald2019}{4}{Calculating IRR does not generate data used in results}. A distinction is drawn between \emph{agreement} (reconciling by discussion) and \emph{reliability} (independent application) \pc{mcdonald2019}{5}{subtle differences between agreement (where two or more coders reconcile differences through discussion)}.

\cf{Critiques of IRR} Alluring to cite for reviewers but possibly harmful: units can alter meaning, some data are interpretable only by the observer, and the pursuit of reliability can reduce sensitivity and marginalise perspectives (feminist HCI) \pc{mcdonald2019}{5}{aspiring to achieve reliability might reduce sensitivity to complex concepts and nuances in data}. Thresholds such as Landis \& Koch are ``clearly arbitrary''. Krippendorff's customary $\alpha\ge.800$, and $\ge.667$ as the lowest limit for tentative conclusions \pc{mcdonald2019}{6}{it is customary to require α ≥ .800}.

\cf{Study design} Sample: random 250 of 617 CSCW and 400 of 1,811 CHI papers (2016--2018), sizes from Cochran's formula (237/317 minimum at 95\,\%/5\,\%). Inclusion: qualitative methods as the primary analysis method; papers using machine learning on text were excluded since the criterion was \emph{methods}, not data \pc{mcdonald2019}{7}{We used Cochran's sample size formula to calculate minimum sample sizes of 237 and 317}. Inclusion coding was checked by two authors: 100\,\% agreement on the first 30, then Cohen's $\kappa=0.849$ and Krippendorff's $\alpha=0.85$ on the next 70, with disagreements discussed and the remainder split \pc{mcdonald2019}{7}{Both of these methods yielded an acceptable level of agreement}. Result: 261 papers (140 CHI, 121 CSCW). Content coding: presence of statements on reliability, number of coders, IRR use and statistic, methodology terms, methods citations, tools. Only unambiguously stated practices were coded. IRR was \emph{not} computed for this coding because it was a binary presence check requiring little interpretation \pc{mcdonald2019}{8}{Because coding involved a simple binary identification}. Excerpts were lightly disguised when critical.

\cf{Findings} Of 251 papers describing analysis: 55.4\,\% described some agreement process, 12.7\,\% used IRR, 48.2\,\% specified the number of researchers (most often two). Cohen's $\kappa$ was the most common statistic (17 of 32), with Fleiss, Krippendorff and Scott less common \pc{mcdonald2019}{10}{Cohen’s kappa was the most commonly used measure of IRR (17 out of 32 papers, 53.1\%)}. Papers rarely justified the statistic. Low values were reported with rationales (e.g.\ $\kappa=0.54$ read as a codebook that was meaningful but not too restrictive) \pc{mcdonald2019}{10}{Some papers conducted a second round of coding and IRR if agreement was too low}. IRR was more frequent for scraped or existing content (18.2\,\%) than for interviews (8.9\,\%) \pc{mcdonald2019}{11}{Table 3: Approaches to reporting on reliability by data collection method}. Papers using a ``codebook'' used IRR in 37\,\% (Table 4). Justifications were sparse \pc{mcdonald2019}{13}{we infer from our readings that many CSCW and HCI researchers assume most readers have sufficient methodological expertise}.

\cf{When NOT to seek agreement or IRR} Codes are the process, not the product; a single expert researcher; social-action or participatory research; agreement can be harmful (feminist and intersectional work); easy or binary coding where one coder suffices; grounded theory, where IRR is not required because codes are an interim product \pc{mcdonald2019}{14}{Even if coders agree on codes, they may interpret the meaning of those codes differently}, \pc{mcdonald2019}{15}{Grounded theory rarely, if ever, requires IRR}.

\cf{When to seek agreement} To code more data across several coders, to confirm measurements, and to identify points of disagreement (and report the data that sparked it) \pc{mcdonald2019}{16}{Identifying where a group of researchers disagrees about interpretations can be helpful}.

\cf{When to use IRR} (1) Consistency between coders coding different subsets. (2) While developing codes. (3) For \emph{predetermined codes}: the more interpretation is required, the more IRR matters \pc{mcdonald2019}{16}{As a general rule with multiple coders, the more interpretation that is required to code the data}. (4) When results are reported quantitatively (frequencies, comparisons) \pc{mcdonald2019}{16}{If findings hinge on frequency counts, there is a strong argument that IRR should be used}. (5) Replicability: state thresholds for each stage, even preregister a target such as 0.7 \pc{mcdonald2019}{17}{Papers could suggest, or even preregister, a target agreement value}. (6) Reducing confirmation bias.

\cf{Reporting advice} Give pragmatic detail on analysis even when obvious. Avoid the passive voice that obscures agency. Avoid ``standard qualitative techniques'' as a shorthand, because no standard exists; describe the process and rationale. Acknowledge resource constraints (few coders, limited expertise) \pc{mcdonald2019}{18}{There is no standard qualitative method in CSCW or HCI}.

\cf{Use for this thesis} A decision aid that frames the walkthrough coding. The thesis applies a \emph{predetermined codebook} to a defined set of observations and reports counts across apps, which is the case in which IRR is recommended (5.3.3--5.3.4), while the inductive steps (codebook development) are the case where discussion matters more. Cite it to justify the two-coder design and to state a priori thresholds, and to explain the reporting items: number of coders, statistic and why, subset size, how disagreements were resolved. Compare with O'Connor \& Joffe for the mechanics, and note that Krippendorff's $\alpha\ge.800$/$.667$ thresholds are quoted here from the 2004 text.
\end{papercard}
